\chapter{Introduction}

\section{Introduction to the Study}

In the digital marketplace, customers continuously express their opinions about brands through online reviews and ratings. For platform-based service businesses such as Zomato, these user-generated reviews are an important source of insight because they reflect real customer experiences related to delivery, refunds, pricing, customer support, and overall app usability. Prior research on electronic word-of-mouth and online reviews shows that such public feedback influences consumer judgment, purchase-related behavior, and brand perception \cite{chevalier2006,filieri2015,park2009}. As the volume of such reviews grows, manually examining them becomes time-consuming and inconsistent, which increases the need for structured analytical methods.

Natural Language Processing (NLP) provides a practical way to organize and interpret large volumes of textual feedback. By applying sentiment analysis and thematic classification to app review data, it becomes possible to identify whether customer opinion is broadly positive, negative, or neutral, and to understand the major issues shaping brand perception. Foundational work in opinion mining and NLP has established sentiment analysis as a useful approach for extracting structured meaning from unstructured text \cite{pang2008,liu2012,manning1999,jurafsky2024}. These insights are valuable not only from a technical standpoint but also from a managerial perspective, as they can support better customer experience decisions and sharper brand strategy.

This study focuses on customer sentiment toward Zomato using Google Play Store reviews. The project uses review ratings as a proxy for sentiment labels and applies rule-based thematic categorization to identify the main drivers of customer opinion. The study generates actionable insights that explain how customers perceive Zomato's service experience and where improvement efforts can create the greatest business value.

\section{Purpose and Scope of the Study}

The purpose of this study is to analyze customer sentiment toward Zomato through Google Play Store reviews and to interpret the results in terms of brand strategy and customer experience. The scope of the study is limited to Zomato-related app reviews collected from the Google Play Store. The core research design and analysis are centered on Zomato, and the final discussion includes a supporting comparison with a key competitor to place the findings in market context.

\section{Research Gap}

Many firms receive a large volume of digital customer feedback, yet decision-making often relies on selective manual reading or surface-level rating summaries. Research on app store feedback has already shown that mobile application reviews are rich in user experience information but difficult to process manually at scale \cite{pagano2013,guzman2017,martin2017}. This creates a gap between available customer voice data and its practical use in managerial decision-making. A focused sentiment and theme-based analysis of Zomato's app reviews can help bridge that gap by converting unstructured customer feedback into interpretable business insights.

\section{Objectives of the Study}

\begin{enumerate}[label=\arabic*.]
    \item To collect and preprocess Zomato-related customer reviews from the Google Play Store.
    \item To classify reviews into positive, negative, and neutral sentiment categories using an NLP-oriented labeling approach.
    \item To identify the major themes influencing customer perception, such as delivery, refunds and cancellations, customer support, pricing and fees, trust, food quality, and app experience.
    \item To examine how review-based sentiment can support brand strategy and customer experience improvement.
    \item To provide managerial insights that can help strengthen Zomato's brand perception and service positioning.
\end{enumerate}

\section{Significance of the Study}

This study is significant from both academic and managerial perspectives. From an academic standpoint, it demonstrates how Natural Language Processing can be applied to large-scale consumer review data to extract structured insights from unstructured textual feedback \cite{pang2008,liu2012}. It also shows how platform-generated ratings and review text can be jointly interpreted to understand customer sentiment in a digital service environment.

From a managerial perspective, the study is relevant because online reviews strongly influence brand image, trust, customer choice, and repeat usage \cite{chevalier2006,filieri2015}. For a platform such as Zomato, understanding what customers appreciate and what they criticize is essential for improving delivery performance, refund handling, app usability, service recovery, and customer communication. The findings of this study therefore support evidence-based decision-making in brand strategy and customer experience management \cite{lemon2016,verhoef2010,keller2013}.

\section{Key Variables of the Study}

The study is centered on a set of core analytical variables derived from Google Play Store review data. These variables are used to interpret customer opinion and identify patterns relevant to business decision-making.

\begin{enumerate}[label=\arabic*.]
    \item \textbf{Sentiment Label:} The overall emotional orientation of a review, classified as positive, negative, or neutral based on the associated star rating.
    \item \textbf{Review Rating:} The numerical score given by the user on the Google Play Store, which serves as the primary proxy for sentiment labeling in this study.
    \item \textbf{Theme Classification:} The dominant issue reflected in a review, such as delivery, refund and cancellation, customer support, pricing and fees, trust, food quality, app experience, order issue, competitor comparison, or other.
    \item \textbf{Review Text:} The unstructured customer feedback used for cleaning, preprocessing, and thematic interpretation.
    \item \textbf{Recency of Review:} The time-based grouping of reviews used to distinguish relatively recent customer opinion from older feedback.
\end{enumerate}

\section{Research Questions}

The study is guided by the following research questions:

\begin{enumerate}[label=\arabic*.]
    \item What is the overall sentiment distribution of Zomato reviews on the Google Play Store?
    \item Which major themes most strongly influence customer perception of Zomato?
    \item What types of issues are most frequently associated with negative customer sentiment?
    \item How can insights derived from app review sentiment support brand strategy and customer experience improvement for Zomato?
\end{enumerate}

\section{Chapter Summary}

This chapter introduced the background, purpose, gap, and objectives of the study. The next chapter will review the relevant literature on sentiment analysis, online customer reviews, NLP-based text analytics, and digital brand perception in service platforms.
